\documentclass[12pt,a4paper]{ctexart}
\usepackage{amsmath,amssymb}
\usepackage{booktabs}
\usepackage{array}
\usepackage{geometry}
\usepackage{enumitem}
\usepackage[hidelinks]{hyperref}
\geometry{margin=2.5cm}
\linespread{1.25}
\setlength{\parskip}{0.4em}
\setlength{\parindent}{2em}
\renewcommand{\arraystretch}{1.15}
\ctexset{section={format={\heiti \zihao{3}}}, subsection={format={\heiti \zihao{4}}}}
\setcounter{secnumdepth}{-1}

\begin{document}

\begin{center}
{\heiti\zihao{3} 2026 CUMCM A 题论文终稿评审报告\par}
\vspace{0.4em}
{\zihao{-4} （评委视角 · 2026-09-13 提交定稿版 · 49 页 · 支撑材料 15 文件）\par}
\end{center}

\section{一、总评与得分}

\textbf{总评：87/100，省一上游，具备冲国奖条件。} 与上一轮模拟评阅（82/100）相比，本轮定稿补齐了附录源程序（8 个程序完整清单，24 页）与支撑材料文件列表，并对 $C_a$ 口径、潜热措辞、守恒表述、结果解释逐条作了收口；经核对，论文附录代码与提交的支撑材料逐字节一致，全部数值（表 5--表 8、$t^*$ 链、Biot 数、7.6 节数量级）内部自洽。当前版本在\textbf{建模严谨性与检验体系}上达到国一水平，主要失分仍在\textbf{图形化呈现与个别表述细节}。

\begin{center}
{\zihao{5}\begin{tabular}{lp{9.5cm}c}
\toprule
分项 & 评语 & 得分 \\
\midrule
摘要 & 方法、结论、检验俱全，问题 1/3/4 有定量结果，问题 2 缺数字 & 9/10 \\
建模与假设 & 四问与题面参数一一对应，假设均有依据与代价分析 & 9/10 \\
求解与计算 & 格式二阶、步长选取流程、四位小数目标明确，结果自洽 & 9.5/10 \\
模型检验 & Bessel 四位一致、收敛比 4.0、守恒闭合、渐近、文献互证，五重检验 & 10/10 \\
灵敏度与稳健性 & 8 变体＋最不利联合工况实算，消融分离收缩与物性 & 9/10 \\
表述与规范 & 图形无引用、符号表不全、双概述图冗余，其余合规 & 7.5/10 \\
\bottomrule
\multicolumn{2}{r}{加权合计} & \textbf{87/100} \\
\bottomrule
\end{tabular}}
\end{center}

\section{二、格式合规核对（2026 规范逐条）}

\begin{center}
{\zihao{5}\begin{tabular}{lp{8.6cm}c}
\toprule
检查项 & 核对结果 & 结论 \\
\midrule
电子版第一页 & 标题＋摘要＋关键词完整在第 1 页 & 通过 \\
页码 & 从摘要页起页脚中部，全文无页眉 & 通过 \\
正文页数 & 一--八节止于第 23 页，$\le$30 页（总 49 页，其中附录代码 24 页） & 通过 \\
匿名 & 无作者、无学校信息字段 & 通过 \\
PDF 大小 & 1.8 MB $\le$ 20 MB & 通过 \\
AI 声明 & 位于参考文献前，含工具、用途、人工核验说明；支撑材料含 AI 工具使用详情.pdf & 通过 \\
附录文件列表 & 列出 result1--4.xlsx、8 个程序、附件副本、AI 详情，与提交的 15 文件逐一对应 & 通过 \\
附录源代码 & 8 个程序完整清单，经逐行比对与支撑材料\textbf{逐字节一致} & 通过 \\
支撑材料 zip & 15 文件、5.3 MB，与文件列表一致 & 通过 \\
\bottomrule
\end{tabular}}
\end{center}

\section{三、主要亮点（评委打高分的依据）}

\begin{enumerate}[label=（\arabic*）,leftmargin=3em]
\item \textbf{四个问题的模型差异逐项对照}（2.4 节表格）：物性经验式、$D$ 的温变性、烘房条件、几何收缩、判据逐列区分，杜绝了参数张冠李戴这一常见失分点。
\item \textbf{数值方法的完整链条}：守恒型有限体积＋Crank--Nicolson＋追赶法，截断误差逐项泰勒展开证明二阶（5.5 节），步长选取以四位小数阈值 $0.5\times10^{-4}$ 为目标、Richardson 实测误差驱动加密（5.6 节）——这是多数国一论文都不具备的严谨性。
\item \textbf{问题 4 的移动边界处理}：从干固体质量守恒与水分守恒律推导随体坐标无对流项形式，不是直接搬公式；再用 2$\times$2 消融把收缩（$-56\%\sim-61\%$）与附录 4 物性差异（$+126\%$）的贡献定量分离，结论"收缩加速占主导、净效应 $-11\%$"有直接数据支撑。
\item \textbf{五重检验全部通过且全部量化}：Bessel 解析解四位小数完全一致；空间收敛比实测 4.0；守恒闭合差 $10^{-19}$（问题 1）至 $10^{-7}$（问题 4）；渐近精确收敛到烘房参数；活化能 32.0 kJ/mol 与三类药材实测值互证。
\item \textbf{不确定度体系自洽}：$t^*=57.14\pm0.03$ h、$50.80\pm0.01$ h 的不确定度来自实测收敛阶外推链，与四位小数精度要求不矛盾；7.6 节对潜热的数量级分析把最大的模型风险以 $+40\%\sim50\%$ 敏感性区间形式显式保留，而非回避。
\end{enumerate}

\section{四、问题清单与定级}

\subsection{4.1 A 类：建议修改（不改不影响成绩，改了更稳）}

\begin{enumerate}[label=A\arabic*.,leftmargin=2.6em]
\item \textbf{全文 7 幅图无一处正文引用。} 没有一句"如图 X 所示"，图与正文脱节，评委阅读时图是孤立插图。建议在 5.4（半径收缩曲线）、6.4（含水率演化、消融图）、7.2（收敛图）、7.7（灵敏度风图）各加一句指引。
\item \textbf{两张概述图功能重叠。} 第 2 页"架构图"与第六节开头"整体解答流程示意图"均为流程型概述图，且图题过于简略（"架构图"三字）。建议二选一，保留的图题改为"本文建模与求解技术路线"，并删去另一张。
\item \textbf{摘要中问题 2 无定量结果。} 其余三问均有数字（36.79 $^\circ$C/1.51 kg/kg、57.14 h、50.80 h），问题 2 只写"给出 0--3 h 的温度与水分场（表 3、表 4）"。建议补一句"3 h 时表面含水率降至 1.01 kg/kg、温度场约 2 h 即与烘房平衡"，摘要四问对称。
\item \textbf{符号表不全。} 正文用到的 $Bi_m$（传质 Biot 数）、$R_g$（气体常数）、$\alpha$（热扩散率）、$\rho_d$（干物质密度）未列入符号说明。建议补 4 行，一分钟的改动，杜绝"符号未定义"类扣分。
\item \textbf{表 1 与表 3 在 0.5 h 的表面温度差异未直接解释。} 同一时刻（1800 s）表 1 表面 36.79 $^\circ$C（附录 2 物性）、表 3 表面 35.41 $^\circ$C（附录 3 物性），差异约 1.4 $^\circ$C，源于 $\rho c_p$、$k$ 不同（附录 3 的 $k$ 更小、$\rho c_p$ 更大，热扩散率低约 14\%）。6.2 节只解释了水分场差异，建议补半句，防评委误认为两表不一致。
\end{enumerate}

\subsection{4.2 B 类：评委质询风险点（论文已防御，答辩需守住）}

\begin{enumerate}[label=B\arabic*.,leftmargin=2.6em]
\item \textbf{潜热未显式建模（假设 H5，全文最大争议点）。} 严格理论下蒸发吸热应作为边界汇项出现，本文以"题给 $D(C,T)$、$h_m$ 为实验标定的有效参数、$E_a=32$ kJ/mol 高于纯水自扩散 18--20 kJ/mol、已含解吸汽化综合能垒"为由不显式叠加，并给出数量级估计：若显式计入，$t^*$ 将增至 79--83 h（$+39\%\sim46\%$），超出题述 2--3 天。防御链完整、态度诚实，但"有效参数已吸收潜热"是观点而非定理，若评委不认同，会认为 $t^*$ 系统性偏低。答辩口径建议：强调"$+40\%\sim50\%$ 已作为敏感性区间计入结果解读，结论的稳健性不受吸收程度判定影响"；切忌声称"潜热一定已完全吸收"。
\item \textbf{$C_a$ 的口径（5.1 节）。} 附件 1 的空气水分单位为 kg/kg（干空气含湿比），本文解释为药材等效平衡水分浓度后与干基含水率 $C$ 同纲相减。物理上两种基准不同，但 7.6 节已作含湿比情景分析、表 8 显示 $C_a\pm20\%$ 对 $t^*$ 影响不足 1\%，且平衡值 0.04986 远低于判据 0.15，风险很低。答辩口径：边界驱动力 $C_R-C_a$ 中 $C_a$ 的贡献始终小于 1\%（表 8 实测）。
\item \textbf{表 7 消融用验证离散，与生产值差 0.05 h。} 附录 3 固定几何值 57.0884 h 与生产值 57.14 h 相差 0.05 h，论文自行声明了这一偏差及理由（同向偏移、百分比结论影响 $<0.1$ 个百分点）。诚实且成立，答辩时照原文口径即可。
\end{enumerate}

\section{五、模拟答辩质询与应对}

\begin{enumerate}[label=Q\arabic*.,leftmargin=2.6em]
\item 问：为什么不建 Luikov 型全耦合模型？答：题给参数体系内信息不足——饱和蒸气压、吸湿等温线均未给出，显式加入潜热汇项需引入题外参数且存在与经验式重复计账的风险；本文以数量级分析（7.6 节）界定其影响上限并计入结论。
\item 问：问题 4 比问题 3 快，收缩效应有这么大吗？答：消融表（表 7）直接回答了这一问题——附录 4 物性单独作用会使时长增至 129 h（$+126\%$），收缩单独作用使时长缩短 56\%--61\%，净效应 $-11\%$；收缩的加速来自 $D/R^2$ 系数随 $R$ 由 2 cm 减至 1.198 cm 增大约 2.8 倍，与 $t^*\propto R^2/D$ 的降速段标度一致。
\item 问：结果保留四位小数，凭什么相信末位？答：5.6 节给出选取流程，7.2 节给出实测——空间步长减半收敛比达理论值 4（问题 1、2），表 5、表 6 的陡梯度行按实测阶 1.3--1.5 作 Richardson 估计，残余误差 $2.7\times10^{-5}\sim3.0\times10^{-5}$，低于四位小数阈值 $0.5\times10^{-4}$ 约半个量级；$t^*$ 经五档加密链外推，离散不确定度仅 0.03 h（问题 3）、0.01 h（问题 4）。
\item 问：烘干时长 57.14 h 与题述 2--3 天是什么关系？答：题面时长是经验范围而非独立实验数据，本文只作数量级一致性参照，不作独立验证；本文的独立验证是解析解、收敛性、守恒、渐近与活化能文献互证五项。
\end{enumerate}

\section{六、结论}

定稿在建模、求解、检验三方面达到国一水准，格式合规核对 9 项全部通过，支撑材料与附录代码逐字节一致、结果可完全复现。剩余问题集中在 A 类表述细节（图引用、符号表、摘要对称性），不影响成绩档位；B 类潜热争议已被论文自身防御并保留敏感性区间，属可答辩的口径问题。\textbf{维持 87/100 的评价：省一上游，具备冲国奖条件}；若答辩前有时间，按 4.1 节五条作最后微调（预计 1 小时内完成），表述分可再上 0.5--1 分。

\end{document}
